\paragraph{Complex matter in a fixed magnetic background.}
The dressed scalar action also admits a controlled complex continuum
limit with a nonzero magnetic field. Here the electromagnetic potential
is prescribed and held fixed in the variation. This is a matter equation
in an external field, rather than an invariant sector of the complete
matter--Maxwell equations.

On the fixed solid \(\Omega\), take
\[
 A(x)=c+\tfrac12 B\times x,\qquad \phi=0,\qquad
 D_A=\nabla-ieA,
\]
where \(c,B\in\mathbb R^3\) and \(e\in\mathbb R\) are fixed.
Use the conforming uniform meshes \(\mathcal T_n\), \(n=2^k\),
of Proposition~\ref{prop:whitney-uniform-cone-refinement}.
For every oriented edge set \(a_{ij}=\int_{v_i}^{v_j}A\cdot dx\).
On a tetrahedron define
\[
 b_i^A(x)=\lambda_i(x)e^{ie\theta_i(x)},\qquad
 \theta_i(x)=\sum_j a_{ij}\lambda_j(x),\qquad
 J_n^Au=\sum_i b_i^A u(v_i).
\]
Let \(V_n^A\) be the complex space obtained from these functions
with one coefficient at each global vertex, including boundary vertices.

\begin{lemma}[Conforming magnetic approximation]
\label{lem:whitney-magnetic-conforming-approximation}
The Whitney interpolant of \(A\) equals \(A\) on every cell, and
\(\theta_i=\int_{v_i}^x A\cdot dx\). The dressed space
\(V_n^A\) is a subspace of \(H^1(\Omega;\mathbb C)\).
There is a constant independent of \(n\) such that
\begin{equation}
 \|J_n^Au-u\|_{H^1}\leq C n^{-1}\|u\|_{H^2}
 \qquad (u\in H^2(\Omega;\mathbb C)).
 \label{eq:whitney-magnetic-interpolation}
\end{equation}
\end{lemma}
\begin{proof}
Write \(A(x)=c+Lx\), where \(L^\mathsf T=-L\).
The local Whitney vector space is exactly
\(\{d+Kx:K^\mathsf T=-K\}\). To check this directly,
\(\lambda_i\nabla\lambda_j-\lambda_j\nabla\lambda_i\)
has skew linear part, and its six oriented edge integrals give the
identity matrix. Thus the six-dimensional space is unisolvent for those
integrals, and reproduces \(A\). This is the lowest-order edge
space in the usual Whitney construction \cite{ArnoldFalkWinther2010}.
Skewness gives
\(\int_{v_i}^x A\cdot dx=A(v_i)\cdot(x-v_i)\),
an affine function of \(x\), whose vertex values are \(a_{ij}\).
On a shared face, the barycentric functions for vertices off that face
vanish. The remaining phases use the same straight paths within the
face, and their traces agree. The resulting piecewise smooth functions
are continuous and therefore belong to \(H^1\). Their vertex values
are the coefficients, so the global basis is linearly independent.

For the approximation bound, fix a cell \(K\) of diameter \(h\)
and a vertex \(x_0\). Put \(A_0=A(x_0)\),
\(\chi(x)=A_0\cdot x\), and \(w=e^{-ie\chi}u\).
Then on that cell
\[
 J_n^Au=e^{ie\chi(x)}
       \sum_i\lambda_i(x)e^{ie\delta_i(x)}w(v_i),
 \quad
 \delta_i=\int_{v_i}^x(A-A_0)\cdot dx.
\]
Uniform shape regularity gives
\(\|\delta_i\|_\infty\leq Ch^2\),
\(\|\nabla\delta_i\|_\infty\leq Ch\), and
\(\|\nabla\lambda_i\|\leq C/h\). Consequently, with
\(r=\sum_i\lambda_i(e^{ie\delta_i}-1)w(v_i)\),
\[
 \|r\|_{L^2(K)}\leq Ch^2 |K|^{1/2}
                         \Bigl(\sum_i|w(v_i)|^2\Bigr)^{1/2},
 \quad
 \|\nabla r\|_{L^2(K)}\leq Ch |K|^{1/2}
                         \Bigl(\sum_i|w(v_i)|^2\Bigr)^{1/2}.
\]
The constants may absorb the fixed upper bound on cell diameters.
Point evaluation on the reference tetrahedron is bounded on \(H^2\)
in three dimensions. Scaling therefore bounds the last weighted nodal
norm by
\(C(\|w\|_{L^2(K)}+h\|\nabla w\|_{L^2(K)}
+h^2\|D^2w\|_{L^2(K)})\).
The ordinary affine interpolation bound for \(I_0w-w\), together
with the estimate for \(r\), gives
\(\|J_n^Au-u\|_{H^1(K)}\leq Ch\|u\|_{H^2(K)}\).
Multiplication by \(e^{\pm ie\chi}\) has uniformly bounded
\(H^2\) operator norm, since \(A_0\) remains bounded on the
fixed solid. Sum the squared local estimates and use \(h<6/n\).
No point evaluation on arbitrary \(H^1\) data is used.
\end{proof}

For \(m^2>0\), use the Hermitian form (linear in its first entry)
\[
 a_A(v,w)=\int_\Omega D_Av\cdot\overline{D_Aw}
                         +m^2v\overline w\,dx.
\]
Its norm is equivalent to \(H^1\), with constants depending only
on \(m^2,e,\|A\|_\infty\). Indeed,
\(\|\nabla v\|_2\leq\|D_Av\|_2+|e|\|A\|_\infty\|v\|_2\),
and the reverse inequality follows by the same triangle inequality.
Define the complex magnetic Ritz projection by
\[
 a_A(R_n^Au,v_n)=a_A(u,v_n)\quad(v_n\in V_n^A).
\]
Orthogonal projection, norm equivalence and
\eqref{eq:whitney-magnetic-interpolation} imply
\begin{equation}
 \|R_n^Au\|_{H^1}\leq C\|u\|_{H^1},\qquad
 \|R_n^Au-u\|_{H^1}\leq Cn^{-1}\|u\|_{H^2}.
 \label{eq:whitney-magnetic-ritz}
\end{equation}

\begin{theorem}[Complex magnetic matter trajectory convergence]
\label{thm:whitney-magnetic-continuum-trajectory}
Fix the preceding background, \(m^2>0\), \(g\geq0\), and a
finite time interval \([0,T]\). Suppose a complex reference
\(u\in C^2([0,T];H^2(\Omega))\) satisfies
\begin{equation}
 (u_{tt},v)+a_A(u,v)+g(|u|^2u,v)=0
                   \quad(v\in H^1(\Omega;\mathbb C)).
 \label{eq:whitney-magnetic-continuum-reference}
\end{equation}
Here \((v,w)=\int v\overline w\). In strong notation this is
\(u_{tt}+D_A^*D_Au+m^2u+g|u|^2u=0\), with magnetic
Neumann condition \(\nu\cdot D_Au=0\). Let \(u_n\in V_n^A\)
satisfy the same equation against all \(v_n\in V_n^A\), with
\(u_n(0)=R_n^Au(0)\), \(\dot u_n(0)=R_n^Au_t(0)\).
Then the finite solution exists for every time and
\begin{equation}
 \sup_{0\leq t\leq T}
 \bigl(\|u_n-u\|_{H^1}+\|\dot u_n-u_t\|_{L^2}\bigr)
       \leq C_T n^{-1}.
 \label{eq:whitney-magnetic-trajectory-error}
\end{equation}
The constant is independent of refinement. These are precisely the
scalar Euler equations of the original dressed action with \(A\)
held fixed, \(\phi=0\), and exact spatial integration.
\end{theorem}
\begin{proof}
The fixed basis has a positive definite mass matrix and no phase velocity.
Taking the real part of the equation tested with \(\dot u_n\)
shows conservation of
\[
 E_n=\tfrac12\|\dot u_n\|_2^2+\tfrac12a_A(u_n,u_n)
                       +\tfrac g4\|u_n\|_4^4.
\]
The initial energies are uniformly bounded by Ritz stability and the
fixed-domain embedding \(H^1\hookrightarrow L^4\).
Coercivity yields uniform \(H^1\) field and \(L^2\) velocity
bounds. At each fixed mesh these also prevent finite-time escape of the
finite polynomial ordinary differential equation, proving global existence.

Put \(\eta=R_n^Au-u\), \(\theta=u_n-R_n^Au\).
The projection is time independent and bounded; it commutes with both
time derivatives. Thus \(\eta,\eta_t,\eta_{tt}\) are
\(O(n^{-1})\) in \(H^1\), uniformly in time. Ritz
orthogonality gives the exact error equation
\[
 (\theta_{tt},v_n)+a_A(\theta,v_n)
 =-(\eta_{tt},v_n)-g(|u_n|^2u_n-|u|^2u,v_n).
\]
For complex numbers the inequality
\(\bigl||z|^2z-|w|^2w\bigr|
\leq 2(|z|^2+|w|^2)|z-w|\), followed by H\"older and
\(H^1\hookrightarrow L^6\), yields
\[
 \||v|^2v-|w|^2w\|_2
 \leq C(\|v\|_{H^1}^2+\|w\|_{H^1}^2)\|v-w\|_{H^1}.
\]
All norms in its coefficient are uniformly bounded by the energy and
the reference hypothesis. Take \(v_n=\dot\theta\), take real
parts, and set
\(\mathcal E=\tfrac12\|\dot\theta\|_2^2+\tfrac12a_A(\theta,\theta)\).
Coercivity and Young's inequality give
\(\dot{\mathcal E}\leq C\mathcal E+Cn^{-2}\).
The initial error energy is zero. Gronwall and the projection estimates
prove the bound. Uniformly energy-bounded alternative initializations
add their initial \(H^1\) field and \(L^2\) velocity distances
from the Ritz data, multiplied by \(C_T\). This proof uses no
inverse inequality or mesh-dependent stability constant.
\end{proof}

The reference hypothesis admits localized complex examples. Extend the
background smoothly with bounded derivatives outside a neighborhood of
the solid, without changing it there. In the magnetic wave equation the
additional connection and potential terms have at most one spatial
derivative. The usual local wave Duhamel contraction in
\(C H^4\cap C^1H^3\), with forcing in \(C H^3\), therefore
applies to smooth compactly supported complex data. The cubic is locally
Lipschitz in this space. The equation gives \(u_{tt}\in C H^2\).
The standard local wave energy argument gives finite propagation speed
one, since the principal part is unchanged \cite{TaoNonlinearWaves}.
For a short interval before the support reaches the boundary, restriction
to the solid has a zero-field boundary collar and obeys the magnetic
Neumann condition. For example, \(u(0)=f\), \(u_t(0)=i\omega f\)
with a nonzero real compact bump and \(e\omega\ne0\) has
\(\operatorname{Im}(\overline u u_t)=\omega f^2\ne0\).
The matter charge need not vanish.

Time independence makes the magnetic mass matrix fixed, and exact Whitney
reproduction removes a potential-consistency defect. Those two facts are
the essential restrictions of this result. In particular, such charged
data generally violate the Gauss equation for \(E=0\); the prescribed
field is not varied, and its sources and backreaction are not derived.
The theorem does not establish the self-consistent charged continuum or
select a physical background, preparation or laboratory clock.

\path{code/electromagnetism/whitney_magnetic_continuum.py} evaluates
nonzero-curvature complex Ritz errors on the same refinements and evolves
a finite nonlinear complex pulse. The independent verifier reconstructs
the mesh, phases and covariant derivatives using a different quadrature
and checks an independent time integration. Gauge changes, shared-face
traces and omitted phase derivatives supply falsifying controls.
These are finite numerical comparisons with declared tolerances;
\eqref{eq:whitney-magnetic-trajectory-error} is an analytic paper proof,
not an interval simulation certificate or a Lean formalization.
The three pointwise norm estimates in
\path{Lean/Screen/MagneticMatterStability.lean} are Lean-checked; they
include the stronger cubic constant \(3/2\). The Sobolev, Ritz and
continuum trajectory arguments above remain analytic proofs.
